\chapter{Operations, limits, and traces}
\index{trace!complex singularity}\index{monotone limit!transcendental metric}
\label{chap:operations-trace}

The inverse-limit definition of
\(\PSH^{\NA}(X,\theta)\) is useful only if the elementary operations on
positive-mass test curves descend through every K\"ahler transition.
They do, but the formulas carry hypotheses: suprema require an upper
bound on terminal parameters, decreasing infima require a lower bound,
and trace exists only on a Lelong-zero domain.  This chapter records
those hypotheses as part of the statements.

In \Crefrange{sec:operations-basic}{sec:operations-volume},
\(X\) is a connected compact K\"ahler manifold, and
\(\theta,\theta'\) are smooth closed forms in pseudoeffective classes.
The operations themselves extend to reduced, irreducible, unibranch
compact K\"ahler spaces with nonempty psh classes by
\cite{XiaOperationsV3}; the smooth setting is retained here so that the
components use \Cref{chap:complex-prerequisites} directly.

The operations are typed.  Addition sends the pair of classes
\((\theta,\theta')\) to \(\theta+\theta'\); translation and maximum stay
in the original class; positive rescaling sends \(\theta\) to
\(\lambda\theta\); and trace replaces \(X\) by the normalization of the
chosen analytic subspace.  These target changes are part of the formulas,
not suppressed identifications.  Likewise, analytic trace outputs are
\(\mathcal I\)-good potentials determined up to \(P\)-equivalence.  The
input dependence is coarser: the resulting \(P\)-class depends only on the
\(\mathcal I\)-class of the ambient potential.  Keeping these two
quotients distinct is essential in the trace formulas below.

\section{Order, addition, translation, and maximum}
\label{sec:operations-basic}

\begin{definition}[Order]
\label{def:operations-order}
For \(\Gamma,\Gamma'\in\PSH^{\NA}(X,\theta)\), write
\(\Gamma\leq\Gamma'\) if, for one equivalently every K\"ahler form
\(\omega\),
\[
 \Gamma_\tau^{\theta+\omega}
 \leq\Gamma_\tau'{}^{\theta+\omega}
 \quad\text{for all }\tau\in\mathbf R.
\tag{54.1}\label{eq:operations-order}
\]
Here both extended curves are used, so the condition also compares
different terminal parameters.
\end{definition}

To compare arbitrary perturbations $\omega_1,\omega_2$, choose
a K\"ahler form $\eta\leq\omega_1,\omega_2$.
If the components are ordered at $\omega_1$, compatibility and
\Cref{lem:tna-transition-order} reflect that order to the
$\eta$-components, then preserve it on projection to $\omega_2$.
This proves the equivalence between comparison at one perturbation
and comparison at every perturbation.
Reflexivity and transitivity are pointwise, and antisymmetry
follows by the same argument in both directions.

\begin{definition}[Addition]
\label{def:operations-addition}
For
\(\Gamma\in\PSH^{\NA}(X,\theta)\) and
\(\Gamma'\in\PSH^{\NA}(X,\theta')\), their sum is the unique element of
\(\PSH^{\NA}(X,\theta+\theta')\) such that
\[
 (\Gamma+\Gamma')^{\theta+\theta'+2\omega}
 =\Gamma^{\theta+\omega}+\Gamma'{}^{\theta'+\omega}
\tag{54.2}\label{eq:operations-addition}
\]
for every K\"ahler form \(\omega\), where the right side is the
test-curve sum \eqref{eq:rays-test-sum}.
\end{definition}

More generally, for positive smooth forms \(\eta,\eta'\),
\[
 P_{\theta+\eta+\theta'+\eta'}
       [\Gamma+\Gamma']_{\mathcal I}
 =
 P_{\theta+\eta}[\Gamma]_{\mathcal I}
 +P_{\theta'+\eta'}[\Gamma']_{\mathcal I}.
\tag{54.3}\label{eq:operations-addition-components}
\]
This formula proves that the symmetric choice \(2\omega\) in
\eqref{eq:operations-addition} does not affect the result.

\begin{definition}[Translation]
\label{def:operations-translation}
For \(C\in\mathbf R\), define \(\Gamma+C\) componentwise by
\[
             (\Gamma+C)^{\theta+\omega}
             :=\Gamma^{\theta+\omega}+C.
\tag{54.4}\label{eq:operations-translation}
\]
On test curves the right side is the parameter translation
\((\Gamma+C)_\tau=\Gamma_{\tau-C}\); on the associated ray it adds
\(Ct\).
\end{definition}

\begin{definition}[Maximum]
\label{def:operations-maximum}
For \(\Gamma,\Gamma'\in\PSH^{\NA}(X,\theta)\), define
\(\Gamma\vee\Gamma'\) by
\[
 (\Gamma\vee\Gamma')^{\theta+\omega}
 :=\Gamma^{\theta+\omega}\vee
       \Gamma'{}^{\theta+\omega}.
\tag{54.5}\label{eq:operations-maximum}
\]
\end{definition}

The right side of \eqref{eq:operations-maximum} is the concave,
$\mathcal I$-model envelope of the two test curves.
The Legendre transform of the raw concave majorant is the pointwise
maximum of the corresponding rays; taking the model envelope can
strictly increase that transform, as in
\Cref{ex:rays-operation-projection}.
Thus \eqref{eq:operations-maximum} uses the projected test-curve
operation of \eqref{eq:rays-test-join}.

\begin{proposition}\label{prop:operations-algebra}
Addition and maximum are commutative and associative.  Translation
satisfies
\[
 \begin{split}
 (\Gamma+\Gamma')+C
   &=(\Gamma+C)+\Gamma'=\Gamma+(\Gamma'+C),\\
 (\Gamma+C)+C'&=\Gamma+(C+C'),\\
 (\Gamma\vee\Gamma')+C
   &=(\Gamma+C)\vee(\Gamma'+C).
 \end{split}
\tag{54.6}\label{eq:operations-algebra}
\]
All identities are equalities of inverse systems, not merely equality
up to \(P\)-equivalence.
\end{proposition}

\begin{proof}
Project to common K\"ahler perturbations, and write
$\alpha,\beta,\gamma$ for the three resulting big backgrounds.
Denote the components again by $\Gamma,\Gamma',\Gamma''$.
Addition is commutative because the substitution
$\delta\mapsto\tau-\delta$ interchanges its two inputs.
For associativity, the model-envelope calculus of
\cite[Proposition~9.4.1]{XiaBook} identifies both bracketings at
each interior parameter with
\[
 P_{\alpha+\beta+\gamma}
 \left[\sup_{\delta,\rho\in\mathbf R}
    \bigl(\Gamma_\delta+\Gamma'_\rho+
                 \Gamma''_{\tau-\delta-\rho}\bigr)\right].
\]
The precise source input here is that inserting the inner model
envelope does not change the final $P$-envelope of this triple
sup-convolution. It does not say that the raw ray sum is geodesic.
The inputs are $\mathcal I$-model, so the same final representative
is obtained with the $\mathcal I$-envelope.

Maximum is commutative and associative by its least-upper-bound
property in \Cref{prop:rays-raw-operations}.
Translation is an order isomorphism with inverse translation by
$-C$, so it commutes with finite joins.
The remaining translation identities follow by replacing
$\tau$ by $\tau-C$ in the defining sup-convolution and then
reindexing its scalar parameter.
The terminal parameters transform by the same sums and maxima.
Finally, \eqref{eq:rays-sum-transition} and maximum compatibility
identify these operations in every perturbation, proving equality
of the compatible inverse systems.
\end{proof}

\begin{example}[Translating the standard action on \(\mathbf P^1\)]
Let \(\Gamma\) be the non-Archimedean metric determined by the action
\(a\cdot[x:y]=[ax:y]\) on
\((\mathbf P^1,\mathcal O(1))\), with the linearization for which
\(x^jy^{m-j}\) has weight \(j\).  The complete calculation is given in
\Cref{ex:rf-p1-complete}; in particular, \(\Gamma_{\max}=1\).

For \(C\in\mathbf Q\), pass to the Veronese degrees for which \(mC\) is an
integer.  Twisting the linearization by \(C\) then changes the degree-\(m\)
weights to \(j+mC\).  Thus the filtration threshold changes by \(mC\), and
the associated test curve and ray satisfy
\[
             (\Gamma+C)_\tau=\Gamma_{\tau-C},
             \qquad
             \ell_t^{\Gamma+C}=\ell_t^\Gamma+Ct,
             \qquad
             (\Gamma+C)_{\max}=1+C.
\]
If \(C,C'\in\mathbf Q\), translating twice changes the weights to
\(j+m(C+C')\), so
\[
                       (\Gamma+C)+C'=\Gamma+(C+C').
\]
The translates $\Gamma+C$ and $\Gamma+C'$ are ordered by their
constants, so their join is the translate with the larger constant.
In this ordered case, model projection leaves the raw ray maximum
unchanged: for $t\geq0$,
\[
 \max\{\ell_t^\Gamma+Ct,\ell_t^\Gamma+C't\}
       =\ell_t^\Gamma+\max\{C,C'\}t,
\]
and hence
\[
             (\Gamma+C)\vee(\Gamma+C')
                =\Gamma+\max\{C,C'\}.
\]
This verifies the translation and maximum identities in
\Cref{prop:operations-algebra} on a nonconstant, explicitly weighted
example.  The formulas for arbitrary real \(C,C'\) follow directly from
the definitions.
\end{example}

\section{Regularized suprema and Choquet reduction}

\begin{definition}[Regularized supremum]
\label{def:operations-supremum}
Let \((\Gamma^i)_{i\in I}\) be a nonempty family in
\(\PSH^{\NA}(X,\theta)\) satisfying
\[
                   \sup_{i\in I}\Gamma^i_{\max}<+\infty.
\tag{54.7}\label{eq:operations-sup-bound}
\]
Its regularized supremum is the unique object with components
\[
 \left(\uscsup_{i\in I}\Gamma^i\right)^{\theta+\omega}
 :=\uscsup_{i\in I}\Gamma^{i,\theta+\omega}.
\tag{54.8}\label{eq:operations-supremum}
\]
\end{definition}

At each component, the right side means the regularized test-curve
supremum: finite maxima are first replaced by their concave
\(\mathcal I\)-model envelopes, then an increasing limit is taken.
Condition \eqref{eq:operations-sup-bound} gives a common upper bound on
the effective domains.

\begin{theorem}\label{thm:operations-choquet}
Under \eqref{eq:operations-sup-bound}, there is a countable
\(I_0\subseteq I\) such that
\[
                   \uscsup_{i\in I}\Gamma^i
                   =\uscsup_{i\in I_0}\Gamma^i.
\tag{54.9}\label{eq:operations-choquet}
\]
\end{theorem}

\begin{proof}
Fix a K\"ahler form \(\omega_0\).  For every \(m\geq1\), apply the
test-curve Choquet lemma in the class
\(\theta+m^{-1}\omega_0\), obtaining a countable subfamily \(I_m\).
Put \(I_0=\bigcup_mI_m\).  If \(\omega\) is arbitrary, choose \(m\)
with \(\omega\geq m^{-1}\omega_0\).  Projecting the equality in the
\(m\)-th component to \(\theta+\omega\), and using commutation of the
\(\mathcal I\)-envelope with regularized suprema, proves equality in
the \(\omega\)-component.  Hence all inverse-system components agree.
\end{proof}

\begin{corollary}\label{cor:operations-increasing-net}
If \((\Gamma^i)_{i\in I}\) is an increasing net satisfying
\eqref{eq:operations-sup-bound}, then its limit is
\(\uscsup_i\Gamma^i\), computed componentwise as in
\eqref{eq:operations-supremum}.
\end{corollary}

\begin{proof}
Every term lies below the regularized supremum.  Conversely, any upper
bound for the increasing net is an upper bound for each finite maximum
and hence for their increasing regularized limit.  Thus the object in
\eqref{eq:operations-supremum} is precisely the order limit, component by
component.
\end{proof}

\begin{remark}[Why nets remain visible]
\label{rem:operations-nets}
The K\"ahler perturbations have the countable cofinal sequence
$(m^{-1}\omega_0)_{m\geq1}$ described in
\Cref{chap:transcendental-metrics}.
This is distinct from the index set $I$ of a family of metrics:
Choquet reduction selects a countable subfamily with the same
supremum, without asserting that its indices are cofinal in $I$.
It also gives no decreasing-limit theorem.
We retain nets in the monotone operation statements for which
the cited source proves net-level continuity.
Trace continuity requires the additional $d_S$-convergence
and domain hypotheses in \Cref{prop:operations-trace-calculus}.
\end{remark}

\clearpage
\section{Decreasing infima and rescaling}

\begin{definition}[Decreasing infimum]
\label{def:operations-infimum}
Let \((\Gamma^i)_{i\in I}\) be a decreasing net satisfying
\[
                   \inf_{i\in I}\Gamma^i_{\max}>-\infty.
\tag{54.10}\label{eq:operations-inf-bound}
\]
Define its infimum by
\[
 \left(\inf_{i\in I}\Gamma^i\right)^{\theta+\omega}
 :=\inf_{i\in I}\Gamma^{i,\theta+\omega}.
\tag{54.11}\label{eq:operations-infimum}
\]
\end{definition}

The existence assertion in \eqref{eq:operations-infimum} requires
positive mass in the limit, not just in each member of the net.
For compatible inverse systems this follows from a smaller
K\"ahler perturbation, as the following proof makes explicit.
The inverse-system operation is
\cite[Definition~4.9 and Lemma~4.10]{XiaOperationsV3}.

\begin{proposition}\label{prop:operations-inf-positive-mass}
Let $(\Gamma^i)_{i\in I}$ be a nonempty decreasing net in
$\PSH^{\NA}(X,\theta)$, and put
$T=\inf_i\Gamma^i_{\max}>-\infty$.
For every K\"ahler form $\omega$ and every $\tau<T$, the potential
$u_\tau^\omega:=\inf_i\Gamma_\tau^{i,\theta+\omega}$
is $\mathcal I$-model and satisfies
\[
 M_{\theta+\omega}(u_\tau^\omega)
       \geq 2^{-n}\int_X\omega^n>0.
\tag{54.11a}\label{eq:operations-inf-mass-bound}
\]
The componentwise infima therefore define an element of
$\PSH^{\NA}(X,\theta)$ with terminal parameter $T$.
It is the greatest lower bound of the net.
\end{proposition}

\begin{proof}
Fix $\omega$ and $\tau<T$, and set
\[
 u_i=\Gamma_\tau^{i,\theta+\omega},
 \qquad
 v_i=\Gamma_\tau^{i,\theta+\omega/2}.
\]
Both nets decrease, and all their members are normalized by
$\sup_X u_i=\sup_X v_i=0$.
The nonempty compact zero sets of either net have the finite
intersection property: a common later index supplies a zero
belonging to any prescribed finite collection.
Consequently neither pointwise infimum collapses to $-\infty$.
The decreasing-net compactness statement in
\Cref{prop:complex-l1-compactness} gives
\[
 u:=\inf_i u_i\in\PSH(X,\theta+\omega),\qquad
 v:=\inf_i v_i\in\PSH(X,\theta+\omega/2).
\]
Both have supremum zero. Envelope monotonicity gives
$u\leq P_{\theta+\omega}[u]_{\mathcal I}\leq u_i$
for every $i$, hence $u=P_{\theta+\omega}[u]_{\mathcal I}$.

Compatibility says
$u_i=P_{\theta+\omega}[v_i]_{\mathcal I}\geq v_i$,
so $u\geq v$.
Put $S=\theta+\omega/2+\ddc v\geq0$.
Mass monotonicity and the nonpluripolar binomial expansion yield
\[
 \begin{split}
 M_{\theta+\omega}(u)
 &\geq M_{\theta+\omega}(v)
   =\int_X\left\langle(S+\omega/2)^n\right\rangle\\
 &\geq\int_X(\omega/2)^n.
 \end{split}
\]
Indeed, every mixed term is positive, and the purely smooth term
retains its full mass because pluripolar sets have smooth volume
zero. This proves \eqref{eq:operations-inf-mass-bound} directly;
it does not assume continuity of mass for arbitrary decreasing
psh potentials.

Concavity is preserved by decreasing limits.
The mass bound now permits application of
\Cref{prop:rays-test-limits} in each fixed background.
The terminal parameter is $T$: every interior level is nontrivial,
whereas for $\tau>T$ some member is identically $-\infty$.
At $T$ use the closed endpoint extension.

For K\"ahler forms $\eta\geq\omega$, the decreasing-envelope
theorem \cite[Proposition~2.8]{XiaOperationsV3}, applied to the
$\mathcal I$-model net $(u_i)_i$ with its now verified positive-mass
limit, gives
\[
 \begin{split}
 P_{\theta+\eta}
       \left[\inf_i\Gamma_\tau^{i,\theta+\omega}\right]_{\mathcal I}
 &=\inf_iP_{\theta+\eta}
       [\Gamma_\tau^{i,\theta+\omega}]_{\mathcal I}\\
 &=\inf_i\Gamma_\tau^{i,\theta+\eta}
 \qquad(\tau<T).
 \end{split}
\tag{54.12}\label{eq:operations-inf-transition}
\]
The endpoint extension makes these equalities compatible as
test curves. Hence the components form an inverse system.
Their pointwise infimum property shows that this system is a
lower bound and dominates every other lower bound.
\end{proof}

The lower bound \eqref{eq:operations-inf-mass-bound} depends
on $\omega$ and tends to zero as the perturbation shrinks.
In \Cref{ex:rays-decreasing-mass-collapse}, the fixed-class masses
tend to zero, precluding compatible smaller components for that net.
Thus inverse systems need only \eqref{eq:operations-inf-bound},
whereas the fixed-big-class theorem retains its mass hypothesis.

\begin{definition}[Positive rescaling]
\label{def:operations-rescaling}
For \(\lambda>0\), define
\(\lambda\Gamma\in\PSH^{\NA}(X,\lambda\theta)\) by
\[
 (\lambda\Gamma)^{\lambda\theta+\omega}
 :=\lambda\Gamma^{\theta+\lambda^{-1}\omega}.
\tag{54.13}\label{eq:operations-rescaling}
\]
\end{definition}

This scales both potentials and the Legendre parameter.  It satisfies
\[
 \begin{split}
 \lambda(\Gamma+\Gamma')&=\lambda\Gamma+\lambda\Gamma',\\
 (\lambda\lambda')\Gamma&=\lambda(\lambda'\Gamma),\\
 \lambda(\Gamma+C)&=\lambda\Gamma+\lambda C,\\
 \lambda\uscsup_i\Gamma^i&=\uscsup_i\lambda\Gamma^i,\\
 \lambda\inf_i\Gamma^i&=\inf_i\lambda\Gamma^i.
 \end{split}
\tag{54.14}\label{eq:operations-rescaling-rules}
\]
No operation with \(\lambda=0\) is defined by
\eqref{eq:operations-rescaling}.

\section{Volume under operations}
\label{sec:operations-volume}

\begin{theorem}\label{thm:operations-volume}
Let the displayed operations be defined.
\begin{enumerate}
\item If \(\Gamma\leq\Gamma'\), then
      \[
                         \vol(\Gamma)\leq\vol(\Gamma').
      \tag{54.15}\label{eq:operations-volume-order}
      \]
\item For \(C\in\mathbf R\),
      \[
                         \vol(\Gamma+C)=\vol(\Gamma).
      \tag{54.16}\label{eq:operations-volume-translation}
      \]
\item For \(\lambda>0\),
      \[
                         \vol(\lambda\Gamma)
                         =\lambda^n\vol(\Gamma).
      \tag{54.17}\label{eq:operations-volume-scale}
      \]
\item If \((\Gamma^i)\) is increasing and satisfies
      \eqref{eq:operations-sup-bound}, then
      \[
      \vol\!\left(\uscsup_i\Gamma^i\right)
                       =\lim_i\vol(\Gamma^i).
      \tag{54.18}\label{eq:operations-volume-increasing}
      \]
\item If \((\Gamma^i)\) is decreasing and satisfies
      \eqref{eq:operations-inf-bound}, only
      \[
      \vol\!\left(\inf_i\Gamma^i\right)
                       \leq\lim_i\vol(\Gamma^i)
      \tag{54.19}\label{eq:operations-volume-decreasing}
      \]
      holds in general.
\end{enumerate}
\par
\end{theorem}

\begin{proof}
For a K\"ahler form $\omega$, abbreviate the initial-level mass by
\[
 V_\omega(\Gamma)
 :=M_{\theta+\omega}
       \bigl(\Gamma_{-\infty}^{\theta+\omega}\bigr).
\]
Thus $\vol(\Gamma)=\inf_\omega V_\omega(\Gamma)$.
Order of test curves implies order of their initial levels, so
mass monotonicity gives \eqref{eq:operations-volume-order}.
Translation leaves the initial level unchanged and proves
\eqref{eq:operations-volume-translation}.
For rescaling, \eqref{eq:operations-rescaling} gives
\[
 M_{\lambda\theta+\omega}
       \bigl((\lambda\Gamma)_{-\infty}^{\lambda\theta+\omega}\bigr)
 =\lambda^n V_{\lambda^{-1}\omega}(\Gamma).
\]
The map $\omega\mapsto\lambda^{-1}\omega$ permutes the K\"ahler
forms, so taking the infimum proves
\eqref{eq:operations-volume-scale}.

We first establish a perturbation bound uniform in the metric.
Fix K\"ahler forms $\omega_0$ and $\Omega$ with
$\Omega\geq\theta+\omega_0$.
For $0<\delta<1$, set
\[
 C(\delta):=
 \sum_{a=0}^{n-1}\binom na
     \int_X\Omega^a\wedge(\delta\omega_0)^{n-a}.
\tag{54.20}\label{eq:operations-uniform-mass-error}
\]
Every summand contains a positive power of $\delta$, hence
$C(\delta)\to0$.
For every $\Gamma\in\PSH^{\NA}(X,\theta)$ we claim
\[
 0\leq V_{\delta\omega_0}(\Gamma)-\vol(\Gamma)
       \leq C(\delta).
\tag{54.20a}\label{eq:operations-uniform-volume-bound}
\]

Write $\omega=\delta\omega_0$, and choose any K\"ahler form
$\omega'\leq\omega$.
Put $\psi=\Gamma_{-\infty}^{\theta+\omega'}$,
$S=\theta+\omega'+\ddc\psi$ and $\gamma=\omega-\omega'$.
The transition mass identity
\Cref{lem:tna-transition-mass-identity} gives
\[
 \begin{split}
 V_\omega(\Gamma)-V_{\omega'}(\Gamma)
 &=\sum_{a=0}^{n-1}\binom na
       \int_X\langle S^a\rangle\wedge\gamma^{n-a}\\
 &\leq\sum_{a=0}^{n-1}\binom na
       \int_X\Omega^a\wedge\omega^{n-a}
   =C(\delta).
 \end{split}
\]
To justify the inequality, first replace $\gamma$ by the larger
semipositive form $\omega$.
The current $R:=\Omega+\ddc\psi$ is positive, with
$R-S=\Omega-\theta-\omega'\geq0$ smooth.
Multilinearity therefore bounds
$\langle S^a\rangle$ by $\langle R^a\rangle$.
Finally, mixed mass monotonicity in the class of $\Omega$
compares the potential $\psi\leq0$ with $0$:
\[
 \int_X\langle R^a\rangle\wedge\omega^{n-a}
       \leq\int_X\Omega^a\wedge\omega^{n-a}.
\]
This is a bound on integrated masses; we do not assert
$S\leq\Omega$ as currents.
The forms $\omega'\leq\omega$ are cofinal among shrinking
K\"ahler perturbations. Taking their infimum proves
\eqref{eq:operations-uniform-volume-bound}, including its
nonnegativity by \Cref{prop:tna-volume-welldefined}.

Now let $(\Gamma^i)_i$ be increasing as in the theorem, and put
$\Gamma=\uscsup_i\Gamma^i$.
At a fixed perturbation, the initial-level identity is
\[
 \Gamma_{-\infty}^{\theta+\omega}
       =\uscsup_i\Gamma_{-\infty}^{i,\theta+\omega}.
\tag{54.20b}\label{eq:operations-increasing-initial-level}
\]
One can check it without interchanging any decreasing limits:
the two regularized suprema, over $i$ and over
$\tau\to-\infty$, are the regularized supremum of the same
family of psh levels.
Equivalently, each side dominates every level
$\Gamma_\tau^{i,\theta+\omega}$ and is bounded by the least
psh upper bound of all these levels.
This is also \cite[Proposition~9.4.3]{XiaBook}.
The initial potentials form an increasing net bounded by zero.
Increasing convergence in $d_S$ and continuity of mass,
\cite[Corollaries~6.2.4--6.2.5]{XiaBook}, give
$V_\omega(\Gamma)=\sup_iV_\omega(\Gamma^i)$.
These are the fixed-background analytic inputs.

Set $L=\sup_i\vol(\Gamma^i)$.
Order already gives $L\leq\vol(\Gamma)$.
By \eqref{eq:operations-uniform-volume-bound}, for every
$0<\delta<1$,
\[
 \vol(\Gamma)
 \leq V_{\delta\omega_0}(\Gamma)
 =\sup_iV_{\delta\omega_0}(\Gamma^i)
 \leq L+C(\delta).
\]
Letting $\delta\to0$ proves
\eqref{eq:operations-volume-increasing}.
The uniform bound justifies interchanging the increasing
metric net with the decreasing perturbation net, as in
\cite[Proposition~13.2.13]{XiaBook}.

For a decreasing net, its infimum lies below each member,
so \eqref{eq:operations-volume-order} gives
\eqref{eq:operations-volume-decreasing}.
The following example shows that the inequality can be strict
even when the limit has positive volume.
\end{proof}

\begin{example}\label{ex:operations-decreasing-volume-drop}
Take $X=\mathbf P^1$ and $\theta=\omega_{\mathrm{FS}}$ with
$\int_X\theta=1$.
Write $\Lambda$ for the hyperplane test curve of
\Cref{ex:rays-hyperplane-test-curve}; thus
$\Lambda_s=0$ for $s\leq0$ and
$M_\theta(\Lambda_s)=1-s$ for $0\leq s\leq1$.
Fix $0<a<1$.
For $j\geq1$ define
\[
 \Gamma^j_\tau=
 \begin{cases}
   \Lambda_{a(1+\tau/j)},&\tau\leq0,\\
   -\infty,&\tau>0.
 \end{cases}
\tag{54.20c}\label{eq:operations-volume-drop-curves}
\]
An increasing affine change of parameter preserves concavity
and decrease, and truncation at the closed endpoint $0$
preserves the test-curve axioms.
Every nontrivial level is $\mathcal I$-model and has mass at least
$1-a>0$.
For fixed $\tau\leq0$, the parameter $a(1+\tau/j)$ increases
with $j$, so the test curves decrease.
They all have terminal parameter zero.

For each $j$, the levels vanish when $\tau\leq-j$, hence
$\Gamma^j_{-\infty}=0$.
At each fixed finite $\tau\leq0$, however, the parameter
increases to $a$, and the closed profile gives
$\Gamma^j_\tau\searrow\Lambda_a$.
The limit is the constant test curve $\Gamma_\tau=\Lambda_a$
for $\tau\leq0$, extended by $-\infty$ for $\tau>0$.
Its initial potential is $\Lambda_a$.
Embedding these positive-mass test curves by
\eqref{eq:tna-positive-embedding}, decreasing-limit compatibility
identifies the same infimum in $\PSH^{\NA}(X,\theta)$.
Formula \eqref{eq:tna-positive-volume-formula} therefore gives
\[
 \vol(\Gamma^j)=1,\qquad
 \vol\!\left(\inf_j\Gamma^j\right)=1-a<1.
\tag{54.20d}\label{eq:operations-strict-volume-drop}
\]

The two limits responsible for the drop are explicit:
for $\tau\leq0$,
\[
 \begin{gathered}
 m_j(\tau):=M_\theta(\Gamma^j_\tau)
       =1-a\max\{1+\tau/j,0\},\\
 \lim_j\lim_{\tau\to-\infty}m_j(\tau)=1,\qquad
 \lim_{\tau\to-\infty}\lim_jm_j(\tau)=1-a.
 \end{gathered}
\]
Thus the obstruction is already the passage to the initial
level in a fixed positive class.
The uniform perturbation estimate does not interchange a
decreasing family with the increasing initial-level supremum.
\end{example}


\section{The analytic trace and its natural domain}

We now return to a smooth compact K\"ahler manifold \(X\) and let
\(Y\subseteq X\) be an irreducible analytic subset.  Let
\(\nu_Y:\widetilde Y\to Y\) be its normalization.  The normalization is
a normal compact K\"ahler space; when a smooth model is needed,
all constructions are made on a resolution and are independent of that
choice.

\begin{definition}[Natural trace domain]
\label{def:operations-trace-domain}
Set
\[
 D_Y
 :=\{\varphi\in\QPSH(X)\mid\nu(\varphi,Y)=0\},
\tag{54.21}\label{eq:operations-trace-domain}
\]
where \(\nu(\varphi,Y)\) is the generic Lelong number along \(Y\).
\end{definition}

The condition in \eqref{eq:operations-trace-domain} is weaker than
\(\varphi|_Y\not\equiv-\infty\).  For a log--log pole, the Lelong number
can vanish while the pointwise restriction to a point is
\(-\infty\).  This is why trace is not naive restriction.

We use the strong singularity pseudometric \(d_S\) of
Darvas--Di~Nezza--Lu.  It is defined on singularity types in a fixed big
class and is complete on its positive-mass part.  Its use below is always
in a fixed background class; no metric is being placed directly on the
set of all quasi-psh functions.

\begin{definition}[Analytic trace]
\label{def:operations-analytic-trace}
Let \(\varphi\in D_Y\), and choose a quasi-equisingular approximation
\((\varphi_j)\) with analytic singularities.  A trace of \(\varphi\)
along \(Y\) is a \(d_S\)-limit
\[
                 \varphi_j|_{\widetilde Y}
                 \xrightarrow{d_S}\psi.
\tag{54.22}\label{eq:operations-analytic-trace}
\]
We write \(\Tr_Y(\varphi)\) for any such \(\psi\).
The output is understood up to \(P\)-equivalence.
\end{definition}

\begin{theorem}\label{thm:operations-trace-well-defined}
For every \(\varphi\in D_Y\), the limit
\eqref{eq:operations-analytic-trace} exists after choosing a sufficiently
positive background class on \(\widetilde Y\).  Every such limit is
\(\mathcal I\)-good, any two such limits are \(P\)-equivalent, and their
common \(P\)-class depends only on the \(\mathcal I\)-class of
\(\varphi\).
\end{theorem}

\begin{proof}[Source boundary and the deductions used here]
Choose a decreasing quasi-equisingular approximation with analytic
singularities.  Vanishing of the generic Lelong number along \(Y\) ensures
that no approximant is identically \(-\infty\) at the generic point of
\(\widetilde Y\).  After adding one sufficiently positive background form,
the restrictions therefore lie in a common positive-mass singularity
space on \(\widetilde Y\).

The deep restriction estimate---that two quasi-equisingular
approximations have restricted tails at \(d_S\)-distance tending to
zero---is the imported input.  It is the trace comparison theorem of
\cite{DX24}, reproduced in the supplied trace chapter of \cite{XiaBook};
its proof uses the characterization of
quasi-equisingular approximation by multiplier ideals and completeness of
the positive-mass \(d_S\)-space.  It makes the restricted sequence Cauchy,
so completeness gives the limit in
\eqref{eq:operations-analytic-trace}, and it gives the same limit
\(d_S\)-class for any second approximation.

Each restricted approximant has analytic singularities.  The
characterization of \(\mathcal I\)-good potentials as the \(d_S\)-closure
of analytic singularity types therefore shows that the limit is
\(\mathcal I\)-good.  Finally, zero \(d_S\)-distance is equivalent to
\(P\)-equivalence, so two limits are \(P\)-equivalent.  If
\(\varphi\sim_{\mathcal I}\varphi'\), the canonical multiplier-ideal
approximants of the two inputs have the same singularity type at every
stage; applying the imported restriction estimate again shows that their
traces have zero \(d_S\)-distance.  This proves the asserted dependence on
the input \(\mathcal I\)-class and yields the finer output \(P\)-class, not
merely the coarser output \(\mathcal I\)-class.
\end{proof}

\begin{proposition}\label{prop:operations-trace-calculus}
Let \(\varphi,\psi\in D_Y\) and \(\lambda>0\).  Whenever the displayed
sum belongs to \(D_Y\),
\[
 \begin{split}
 \varphi\preceq_{\mathcal I}\psi
 &\Longrightarrow
 \Tr_Y(\varphi)\preceq_P\Tr_Y(\psi),\\
 \Tr_Y(\varphi+\psi)
 &\sim_P\Tr_Y(\varphi)+\Tr_Y(\psi),\\
 \Tr_Y(\lambda\varphi)
 &\sim_P\lambda\Tr_Y(\varphi),\\
 \Tr_Y(\varphi\vee\psi)
 &\sim_P\Tr_Y(\varphi)\vee\Tr_Y(\psi).
 \end{split}
\tag{54.23}\label{eq:operations-trace-calculus}
\]
If \((\varphi_j)_{j\in J}\) is a decreasing net in one common psh class,
\(\varphi_j\xrightarrow{d_S}\varphi\), and \(\varphi\in D_Y\), then
\[
                  \Tr_Y(\varphi_j)
                  \xrightarrow{d_S}\Tr_Y(\varphi).
\tag{54.24}\label{eq:operations-trace-decreasing}
\]
\end{proposition}

\begin{proof}[Source boundary and quotient check]
Choose analytic-singularity approximants for \(\varphi\) and \(\psi\) on
one common index sequence.  Restriction commutes with addition, positive
scaling, and finite maximum at this analytic stage; order of singularities
is preserved as well.  Passing to the \(d_S\)-limits gives the first four
relations in \(P\)-order or \(P\)-equivalence, since those are respectively
closed order and equality in the \(d_S\)-space.  Compatible canonical
approximants show that the first relation depends only on
\(\mathcal I\)-order of the inputs.

The decreasing-net assertion is the trace continuity theorem of
\cite{DX24}, reproduced in the supplied trace chapter of \cite{XiaBook}.
Its nonformal input is the same uniform restriction estimate used in
\Cref{thm:operations-trace-well-defined}.  The source first proves the
sequential assertion by diagonal approximation.  If the net statement
failed, a positive \(d_S\)-gap from the alleged limit would allow one to
choose recursively an increasing sequence of indices that retains that
gap.  The corresponding decreasing sequence still converges to
\(\varphi\) in \(d_S\), contradicting the sequential assertion.  This
reduces the net conclusion to the imported sequential estimate.
\end{proof}

\begin{remark}[No increasing-continuity claim]
\label{rem:operations-trace-not-increasing}
Trace is not continuous along arbitrary increasing sequences.  Locally,
\[
 \varphi_j(z,w)
 =\log|z|^2\vee j^{-1}\log|w|^2+O(1)
\]
increases almost everywhere to a bounded singularity, while on
\(Y=\{w=0\}\) every restriction retains the pole \(\log|z|^2\).
\end{remark}

This example also explains why zero generic Lelong number, rather than
pointwise finiteness on \(Y\), is the natural domain condition.  A local
log--log pole can equal \(-\infty\) on \(Y\) while its generic Lelong
number is zero; analytic trace is defined by equisingular approximation
and still makes sense.  Naive substitution on \(Y\) would discard that
case.

\section{Descent to multiplier-ideal classes}

The quotient in the next statement is the one used in the source.

\begin{theorem}\label{thm:operations-trace-p-descent}
The assignment
\[
 [\varphi]_{\mathcal I}\longmapsto
 [\Tr_Y(\varphi)]_P
\tag{54.25}\label{eq:operations-trace-p-map}
\]
defines a map
\[
 D_Y/{\sim_{\mathcal I}}
 \longrightarrow
 \QPSH(\widetilde Y)/{\sim_P}.
\tag{54.26}\label{eq:operations-trace-p-descent}
\]
\end{theorem}

\begin{proof}
First, \(D_Y\) is saturated for \(\mathcal I\)-equivalence.  Indeed,
\(\mathcal I\)-equivalence is equality of Lelong numbers on every
modification.  Blowing up the generic point of \(Y\) recovers the
generic Lelong number along \(Y\), so
\[
 \varphi\sim_{\mathcal I}\psi
 \quad\Longrightarrow\quad
 \nu(\varphi,Y)=\nu(\psi,Y).
\tag{54.27}\label{eq:operations-domain-saturated}
\]
Applying the first line of \eqref{eq:operations-trace-calculus} in both
directions shows that \(\mathcal I\)-equivalent inputs have
\(P\)-equivalent traces.  This proves
\eqref{eq:operations-trace-p-descent}; it is also the dependence assertion
of \cite{DX24}, reproduced in the supplied trace chapter of
\cite{XiaBook}.
\end{proof}

\begin{remark}[Two different quotients]
\label{rem:operations-trace-p-boundary}
\Cref{thm:operations-trace-p-descent} has an
\(\mathcal I\)-class as input and a \(P\)-class as output.  Since
\(P\)-equivalence implies \(\mathcal I\)-equivalence, it also yields a map
from ambient \(P\)-classes, but that is a weaker input statement.  Since
output \(P\)-equivalence implies output \(\mathcal I\)-equivalence, one may
also forget to the coarser output quotient.  The content of the theorem is
the stronger cross-quotient map displayed in
\eqref{eq:operations-trace-p-descent}; it does not choose a distinguished
potential inside the output \(P\)-class.
\end{remark}

Because every analytic trace is \(\mathcal I\)-good,
\[
 P_\beta[\Tr_Y(\varphi)]
 =P_\beta[\Tr_Y(\varphi)]_{\mathcal I}
\]
whenever either side is defined in the positive-mass
\(\beta\)-psh space.  Thus the two envelope notations agree on trace
outputs even though they do not agree on arbitrary quasi-psh functions.

\section{Trace of a transcendental metric}

\begin{definition}[Natural domain for NA trace]
\label{def:operations-na-trace-domain}
Let \(\Gamma\in\PSH^{\NA}(X,\theta)\).  Its trace along \(Y\) is defined
when
\[
 \nu(\Gamma_\tau^{\theta+\omega},Y)=0
\tag{54.28}\label{eq:operations-na-trace-domain}
\]
for every K\"ahler form \(\omega\) and all sufficiently negative
\(\tau\).  In that case put
\[
 (\Tr_Y\Gamma)_{\max}
 :=\sup\{\tau<\Gamma_{\max}\mid
          \nu(\Gamma_\tau^{\theta+\omega},Y)=0\}.
\tag{54.29}\label{eq:operations-na-trace-max}
\]
This value is independent of \(\omega\).
\end{definition}

For a K\"ahler form \(\eta\) on \(\widetilde Y\), choose a K\"ahler
form \(\widetilde\omega\) on \(X\) with
\(\eta\geq\widetilde\omega|_{\widetilde Y}\).  For
\(\tau<(\Tr_Y\Gamma)_{\max}\), set
\[
 (\Tr_Y\Gamma)_\tau^{\theta|_{\widetilde Y}+\eta}
 :=
 P_{\theta|_{\widetilde Y}+\eta}
 \left[
 \Tr_Y^{\theta+\widetilde\omega}
   \bigl(\Gamma_\tau^{\theta+\widetilde\omega}\bigr)
 \right]_{\mathcal I}.
\tag{54.30}\label{eq:operations-na-trace}
\]

\begin{proposition}\label{prop:operations-na-trace}
Formula \eqref{eq:operations-na-trace} is independent of all choices and
defines
\[
 \Tr_Y\Gamma\in
 \PSH^{\NA}(\widetilde Y,\theta|_{\widetilde Y}).
\tag{54.31}\label{eq:operations-na-trace-target}
\]
If \(\pi:X'\to X\) is proper bimeromorphic, \(Y\) is not contained in
its non-isomorphism locus, and \(Y'\) is the strict transform, then
\[
 \widetilde\pi^*\Tr_Y\Gamma
 =\Tr_{Y'}(\pi^*\Gamma)
\tag{54.32}\label{eq:operations-na-trace-birational}
\]
on the normalizations.
\end{proposition}

\begin{proof}
If two ambient perturbations are admissible in
\eqref{eq:operations-na-trace}, move both to a common larger perturbation.
Compatibility of the DXZ components makes the input levels
\(\mathcal I\)-equivalent.  By
\Cref{thm:operations-trace-well-defined}, their analytic traces are
\(P\)-equivalent, hence also \(\mathcal I\)-equivalent; taking the
displayed \(\mathcal I\)-envelope in one
target class makes the normalized levels equal.  Repeating the argument
for two target perturbations proves the inverse-system cocycle.  Order and
concavity follow from \Cref{prop:operations-trace-calculus}, and the domain
condition ensures the required positive-mass interior levels after target
perturbation.  Since analytic traces are \(\mathcal I\)-good, the displayed
\(\mathcal I\)-envelopes agree with their \(P\)-envelopes.  Hence the
levels form the metric of
\eqref{eq:operations-na-trace-target}.  This is
\cite[Definition~4.20]{XiaOperationsV3}.

For \eqref{eq:operations-na-trace-birational}, pull back
quasi-equisingular approximants.  Away from the non-isomorphism locus their
restrictions to the strict transform agree with the pullbacks of the
original restrictions.  Passing first to \(d_S\)-limits and then to the
normalized envelopes proves the identity level by level, as in
\cite[Section~4.8]{XiaOperationsV3}.
\end{proof}

\begin{remark}[Trace domain firewall]
\label{rem:operations-trace-firewall}
Condition \eqref{eq:operations-na-trace-domain} is part of the
definition.  It is not enough to check one terminal level, nor may one
replace it by nontrivial pointwise restriction.  The supremum in
\eqref{eq:operations-na-trace-max} may be strictly smaller than
\(\Gamma_{\max}\), and levels at or above it are governed by the
extended convention rather than by analytic restriction.
\end{remark}

\begin{remark}[Source note]
\label{rem:operations-source-note}
The componentwise operations, Choquet theorem, monotone-limit rules,
volume identities, and transcendental trace follow the operations
sections of the supplied lectures and \cite{XiaOperationsV3}.
Numbered references to the operations paper in this chapter use
arXiv version~3, dated 22 October 2024; the published chapter
\cite{Xia23Operations} is a separate bibliography entry.
The order-reflection input is also
\cite[Proposition~6.1.3]{XiaBook}.
The mass bound in \Cref{prop:operations-inf-positive-mass}
verifies the positive-mass prerequisite of the decreasing-envelope
theorem before its application to inverse systems.  The analytic trace, its comparison
calculus, and its continuity follow \cite{DX24} and the complete supplied
trace chapter.  The hypotheses
used there are retained: broad operations extend to reduced,
irreducible, unibranch compact K\"ahler spaces, whereas the analytic
trace and \Cref{thm:operations-trace-p-descent} are stated here on the
smooth ambient manifold and on the natural domain \(D_Y\).  In
particular, the supplied trace chapter proves that every trace is
\(\mathcal I\)-good, that two choices are \(P\)-equivalent, and that the
resulting \(P\)-class depends only on the ambient \(\mathcal I\)-class.  Its
continuity theorem is stated for decreasing nets under the common-class,
\(d_S\)-convergence, and zero-generic-Lelong-number hypotheses recorded in
\Cref{prop:operations-trace-calculus}.  These source-level distinctions are
not weakened here to an \(\mathcal I\)-valued trace or strengthened beyond
the stated net theorem.
\end{remark}
